% from ega3-1-fr.tex
\oldpage[0\textsubscript{III}]{34}
$\mathrm R^pT$ désignant comme d'ordinaire le $p$-ième \emph{foncteur
dérivé} de $T$ pour $p\in\mathbb Z$~; $\mathrm R^qT(K^\bullet)$ désigne le
complexe $(\mathrm R^qT(K^i))_{i\in\mathbb Z}$.

\phantomsection\label{0.11.4.4-fr}
Sauf mention expresse du contraire, nous supposerons désormais que tout objet
de $\mathcal C$ est sous-objet d'un objet injectif de $\mathcal C$, de sorte
que les résolutions de Cartan-Eilenberg injectives existent pour tout complexe
de $\mathcal C$. Comme $L^{ij}=0$ pour $j<0$, les critères de
(\hyperref[0.11.3.3-fr]{11.3.3}) montrent que les \emph{deux} suites
spectrales d'hypercohomologie de $T$ par rapport à $K^\bullet$ existent et
sont \emph{birégulières} dans chacun des deux cas suivants~:
\begin{enumerate}
  \item[$1^\circ$] $K^\bullet$ est limité inférieurement~;
  \item[$2^\circ$] Tout objet de $\mathcal C$ admet une résolution injective
    de longueur au plus égale à un entier $n$ (indépendant de l'objet
    considéré).
\end{enumerate}
En effet, dans le premier cas, on peut supposer
(\hyperref[0.11.4.2-fr]{11.4.2}) qu'il existe $i_0$ tel que $L^{ij}=0$ pour
$i<i_0$ et dans le second qu'il existe $j_1$ tel que $L^{ij}=0$ pour $j>j_1$~;
dans chacun des deux cas, il est clair en outre que pour $n$ donné, il n'y a
qu'un nombre fini de couples $(i,j)$ tels que $L^{ij}\neq0$ et $i+j=n$, ce
qui établit nos assertions.

Lorsqu'on suppose que dans $\mathcal C'$ les limites inductives filtrantes
existent et sont exactes (ce qui implique en particulier l'existence dans
$\mathcal C'$ des sommes directes infinies), alors le complexe défini par le
bicomplexe $T(L^{\bullet\bullet})$ existe, et le critère
(\hyperref[0.11.3.3-fr]{11.3.3}) montre que la suite
$'E(T(L^{\bullet\bullet}))$ est toujours \emph{régulière}.

Lorsque $K^\bullet$ est un complexe dont tous les termes $K^i$ sont nuls sauf
un seul $K^{i_0}$, $\mathrm R^nT(K^\bullet)$ est isomorphe à
$\mathrm R^{n-i_0}T(K^{i_0})$, ainsi qu'il résulte aussitôt des définitions en
prenant une résolution de Cartan-Eilenberg $L^{\bullet\bullet}$ telle que
$L^{ij}=0$ pour $i\neq i_0$.

Si $K^\bullet$ et $K'^\bullet$ sont deux complexes de $\mathcal C$, $f$, $g$
deux morphismes \emph{homotopes} de $K^\bullet$ dans $K'^\bullet$, alors les
morphismes $\mathrm R^\bullet T(K^\bullet)\to
\mathrm R^\bullet T(K'^\bullet)$ déduits de $f$ et $g$ sont identiques, et il
en est de même pour les morphismes des suites spectrales de cohomologie.
\end{env}

\begin{proposition}[11.4.5]
\label{0.11.4.5-fr}
Supposons que dans $\mathcal C'$ les limites inductives filtrantes existent et
soient exactes. Si $\mathrm R^nT(K^i)=0$ pour tout $n>0$ et tout
$i\in\mathbb Z$, on a des isomorphismes fonctoriels
\[
  \label{0.11.4.5.1-fr}
  \mathrm R^iT(K^\bullet)\xrightarrow{\sim}
  \mathrm H^i(T(K^\bullet))\qquad(i\in\mathbb Z).
  \tag{11.4.5.1}
\]
\end{proposition}

\begin{proof}
En effet, les seuls termes $E_2$ non nuls de la première suite spectrale
(\hyperref[0.11.4.3.1-fr]{11.4.3.1}) sont alors
$'E_2^{p0}=\mathrm H^p(T(K^\bullet))$~; autrement dit, cette suite est
\emph{dégénérée}~; comme elle est régulière
(\hyperref[0.11.4.4-fr]{11.4.4}), la conclusion résulte de
(\hyperref[0.11.1.6-fr]{11.1.6}).
\end{proof}

\begin{env}[11.4.6]
\label{0.11.4.6-fr}
Considérons maintenant, par exemple, un \emph{bifoncteur covariant}
$(M,N)\mapsto T(M,N)$ de $\mathcal C\times\mathcal C'$ dans $\mathcal C''$,
$\mathcal C$, $\mathcal C'$, $\mathcal C''$ étant trois catégories
abéliennes~; on suppose, pour simplifier, $T$ additif en chacun de ses
arguments, et en outre que tout objet de $\mathcal C$ et tout objet de
$\mathcal C'$ sont sous-objets d'un objet injectif, et que les limites
inductives filtrantes existent dans $\mathcal C''$ et sont exactes. On définit
alors l'\emph{hypercohomologie de $T$ par rapport à deux complexes}
$K^\bullet$, $K'^\bullet$ de $\mathcal C$ et $\mathcal C'$ respectivement, à
opérateurs de dérivation de degré $+1$, en prenant pour $K^\bullet$ (resp.
$K'^\bullet$) une résolution de Cartan-Eilenberg injective
$L^{\bullet\bullet}$ (resp. $L'^{\bullet\bullet}$)~; alors
$T(L^{\bullet\bullet},L'^{\bullet\bullet})$ est un quadricomplexe de
$\mathcal C''$, que l'on considère comme \emph{bicomplexe} de $\mathcal C''$
en prenant pour degrés de $T(L^{ij},L'^{hk})$ les entiers $i+h$ et $j+k$.
L'\emph{hypercohomologie de $T$ par rapport à $K^\bullet$ et $K'^\bullet$}
est par définition la cohomologie
$\mathrm H^\bullet(T(L^{\bullet\bullet},L'^{\bullet\bullet}))$ de ce
bicomplexe (autrement dit, celle du complexe simple associé) notée
$\mathrm R^\bullet T(K^\bullet,K'^\bullet)$~; elle est l'aboutissement de
deux suites spectrales dont les termes $E_2$ sont donnés par
